% =============================================================================
% CHAPTER 1: OVERVIEW AND TRACEABILITY METHODOLOGY
% IP Traceability Specification
% =============================================================================

\chapter{Overview and Traceability Methodology}
\label{chap:overview}

\section{Purpose}

This specification provides \textbf{auditable evidence} demonstrating how the intellectual property (IP) described in nine source documents has been implemented---through conceptual adaptation---into the SAP-OSS codebase. The purpose is not to show code copying, but to evidence that the \textit{ideas, concepts, and architectural patterns} from the IP papers have shaped the structures and implementations in the codebase.

\section{Scope}

\subsection{IP Sources Traced}

This specification traces concepts from the following nine IP documents to their implementations:

\begin{table}[H]
\centering
\caption{IP Source Documents and Primary Implementation Targets}
\label{tab:ip-sources}
\begin{tabularx}{\textwidth}{@{}llX@{}}
\toprule
\textbf{IP ID} & \textbf{Document} & \textbf{Primary Codebase Target} \\
\midrule
IP-001 & safety-ai-governance & \filepath{src/intelligence/ai-core-pal/governance/} \\
IP-002 & safety-ndeductivedatabase-mangle & Datalog/governance concepts \\
IP-003 & safety-nlocalmodels-mojo & \filepath{src/intelligence/vllm-turboquant/} \\
IP-004 & safety-ntimeseries-mojo & \filepath{src/intelligence/analytics/} \\
IP-005 & safety-aimo-ntimeseries & Mathematical foundations \\
IP-006 & AI Earnings Analysis Module & NLP/analytics concepts \\
IP-007 & Ai Surprise\_small & \filepath{src/intelligence/analytics/anomaly\_detector.py} \\
IP-008 & digitaltwin & World model concepts \\
IP-009 & narrariveai & Stability/governance concepts \\
\bottomrule
\end{tabularx}
\end{table}

\subsection{Codebase Scope}

The implementation evidence is drawn from the following codebase directories:

\begin{itemize}[nosep]
    \item \filepath{src/intelligence/} -- AI core services, analytics, OCR, vLLM
    \item \filepath{src/training/} -- Training pipelines and schema integration
    \item \filepath{src/generativeUI/} -- Generative UI components
    \item \filepath{src/data/} -- LangChain integration, OData vocabularies
\end{itemize}

\section{Traceability Methodology}

\subsection{Evidence Types}

Each IP concept is traced using the following evidence hierarchy:

\begin{enumerate}[nosep]
    \item \textbf{Verbatim Quotation}: Exact text from the IP source document
    \item \textbf{Concept Identification}: The specific idea or pattern being traced
    \item \textbf{Implementation Location}: File path, class, or function name
    \item \textbf{Code Evidence}: Actual code snippet demonstrating the concept
    \item \textbf{Adaptation Notes}: How the concept was adapted (not copied)
\end{enumerate}

\subsection{Traceability Levels}

Implementation evidence is provided at three levels of granularity:

\begin{table}[H]
\centering
\caption{Traceability Granularity Levels}
\label{tab:trace-levels}
\begin{tabularx}{\textwidth}{@{}lXl@{}}
\toprule
\textbf{Level} & \textbf{Description} & \textbf{Example} \\
\midrule
Service & High-level architectural component & \filepath{src/intelligence/} \\
Module & Specific Python/Mojo module & \filepath{analytics/anomaly\_detector.py} \\
Function & Individual function or method & \funcname{\_mad\_detection()} \\
\bottomrule
\end{tabularx}
\end{table}

\subsection{Evidence Format}

Each traced concept follows this standardized format:

\begin{ipsourcebox}{IP-XXX, Section Y}
\textit{``Verbatim quoted text from the IP document describing the concept...''}
\end{ipsourcebox}

\begin{tracerecord}
\textbf{IP Source ID} & IP-XXX \\
\textbf{IP Document} & Document filename \\
\textbf{IP Location} & Section, Page, Lines \\
\textbf{Concept} & Brief concept description \\
\midrule
\textbf{Implementation Path} & \filepath{src/path/to/file.py} \\
\textbf{Key Files} & List of implementing files \\
\textbf{Functions/Classes} & \funcname{ClassName}, \funcname{function\_name} \\
\textbf{Adaptation Type} & Conceptual / Algorithmic / Structural \\
\end{tracerecord}

\begin{implbox}{Implementation Evidence}
Code snippet demonstrating the concept implementation.
\end{implbox}

\begin{adaptbox}
Explanation of how the IP concept was adapted into the implementation, noting any differences between the original concept and its realization.
\end{adaptbox}

\section{IP to Codebase Mapping Overview}

\subsection{Conceptual Flow}

The following diagram illustrates how IP concepts flow into the codebase:

\begin{center}
\begin{tikzpicture}[
    node distance=1.5cm,
    ipbox/.style={rectangle, draw=sourcegreen, fill=sourcegreen!10, 
                  text width=3cm, align=center, rounded corners, minimum height=1cm},
    codebox/.style={rectangle, draw=implblue, fill=implblue!10,
                    text width=3cm, align=center, rounded corners, minimum height=1cm},
    arrow/.style={->, >=stealth, thick}
]
    % IP Documents
    \node[ipbox] (ip1) {IP-001--IP-005\\Safety Papers};
    \node[ipbox, below=of ip1] (ip2) {IP-006--IP-009\\AI/ML Papers};
    
    % Adaptation Layer
    \node[rectangle, draw=orange, fill=orange!10, 
          text width=2.5cm, align=center, rounded corners,
          right=2cm of ip1, yshift=-0.75cm] (adapt) {Concept\\Adaptation};
    
    % Codebase
    \node[codebox, right=2cm of adapt, yshift=0.75cm] (code1) {src/intelligence/\\governance/};
    \node[codebox, below=of code1] (code2) {src/intelligence/\\analytics/};
    
    % Arrows
    \draw[arrow] (ip1.east) -- (adapt.west |- ip1.east);
    \draw[arrow] (ip2.east) -- (adapt.west |- ip2.east);
    \draw[arrow] (adapt.east |- code1.west) -- (code1.west);
    \draw[arrow] (adapt.east |- code2.west) -- (code2.west);
\end{tikzpicture}
\end{center}

\subsection{Key Concept Categories}

The IP documents contribute concepts in four primary categories:

\begin{table}[H]
\centering
\caption{IP Concept Categories}
\label{tab:concept-categories}
\begin{tabularx}{\textwidth}{@{}lXl@{}}
\toprule
\textbf{Category} & \textbf{Description} & \textbf{IP Sources} \\
\midrule
Safety/Governance & Blanket control, audit trails, safety gates & IP-001, IP-002 \\
Numerical Stability & Softmax, RoPE, quantization bounds & IP-003, IP-005 \\
Analytics/Anomaly & MAD, Z-score, Kalman, FFT & IP-004, IP-007 \\
AI Architecture & LNN, world models, stability metrics & IP-006, IP-008, IP-009 \\
\bottomrule
\end{tabularx}
\end{table}

\section{Document Structure}

The remainder of this specification is organized as follows:

\begin{itemize}[nosep]
    \item \textbf{Chapters 2--10}: Individual IP document traceability (one chapter per IP)
    \item \textbf{Chapter 11}: Consolidated traceability matrix
    \item \textbf{Chapter 12}: Concept adaptation registry
    \item \textbf{Appendix A1}: Complete IP source document index
    \item \textbf{Appendix A2}: Code reference guide
\end{itemize}

Each chapter follows the same structure:
\begin{enumerate}[nosep]
    \item IP Document Summary
    \item Key Concepts Identified
    \item Traceability Evidence (per concept)
    \item Implementation Summary
\end{enumerate}